
\documentclass[../../../physics-standard_answers_all.tex]{subfiles}

\begin{document}

\setlength{\abovedisplayskip}{2mm}
\setlength{\belowdisplayskip}{1mm}

\begin{enumerate}

    \item
        $ \displaystyle \hspace{1\zw}
        \begin{aligned}[t]
            E_x(x,\hspace{1.2mm} 0)
            & = k_0 \frac{Q}{L^2+x^2} \cdot \frac{x}{\sqrt{L^2+x^2}} \\
            & \quad\: + k_0\frac{Q}{L^2+x^2} \cdot \frac{x}{\sqrt{L^2+x^2}} \\
            & = \uwave{\frac{2k_0Qx}{\left(L^2+x^2\right)^{3/2}}} \;.
        \end{aligned}
        $
        \begin{align*}
            & E_y(x,\hspace{1.2mm} 0)
            \begin{aligned}[t]
                & = - k_0 \frac{Q}{L^2+x^2} \cdot \frac{L}{\sqrt{L^2+x^2}} \\
                & \quad\: + k_0 \frac{Q}{L^2+x^2} \cdot \frac{L}{\sqrt{L^2+x^2}} \\
                & = \uwave{0}\;. 
            \end{aligned} \\
            & \phi(x,\hspace{1.2mm}0) = k_0\frac{Q}{\sqrt{L^2+x^2}}
            + k_0\frac{Q}{\sqrt{L^2+x^2}} = \uwave{\frac{2k_0Q}{\sqrt{L^2+x^2}}} \;.
        \end{align*}

    \item 仕事をエネルギー収支から逆算する．
        \begin{align*}
            W = q \phi (\sqrt{3}L,\hspace{1.2mm}0) - q \cdot 0 = \uwave{\frac{k_0 Qq}{L}}\;.
        \end{align*}

    \item 位置エネルギーが最大（運動エネルギーが最小）となる原点Oを通過できればよい．
        原点Oに到達したときの速さを$u$として，力学的エネルギー保存則より，
        %
        \begin{align*}
            \frac{1}{2}mu^2 + q \phi(0,\hspace{1.2mm}0)
            = \frac{1}{2}m{v_0}^2 + q \phi(\sqrt{3}L,\hspace{1.2mm}0) \;. 
        \end{align*}
        %
        よって，$ \displaystyle \frac{1}{2}mu^2 = \frac{1}{2}m{v_0}^2
        - \frac{k_0 Qq}{L} > 0$となるために，
        %
        \begin{align*}
            \uwave{v_0 > \sqrt{\frac{2k_0 Qq}{mL}}} \;.
        \end{align*}

    \item 仕事をエネルギー収支から逆算する．
        \begin{align*}
            W^* = (-q) \phi (\sqrt{3}L,\hspace{1.2mm}0) - (-q) \cdot 0
            = \uwave{- \frac{k_0 Qq}{L}}\;.
        \end{align*}

    \item 無限遠に到達したときの速さを$v_\infty$として，力学的エネルギー保存則より，
        %
        \begin{align*}
            \frac{1}{2}m{v_\infty}^2 + (-q) \cdot 0
            = \frac{1}{2}m{v_0}^2 + (-q) \phi(\sqrt{3}L,\hspace{1.2mm}0) \;. 
        \end{align*}
        %
        よって，
        $ \displaystyle    
            \frac{1}{2}m{v_\infty}^2
            = \frac{1}{2}m{v_0}^2 - \frac{k_0 Qq}{L} > 0
        $
        となるために，
        %
        \begin{align*}
            \uwave{v_0 > \sqrt{\frac{2k_0 Qq}{mL}}} \;.
        \end{align*}

    \item 運動方程式（$x$方向）より，
        %
        \begin{align*}
            & m a_x = (-q) E_x(x,\hspace{1.2mm}0)
            = -\frac{2k_0 Qq}{\left(L^2+x^2\right)^{3/2}} x \\
            & \text{∴} \;\; a_x \fallingdotseq -\frac{2k_0 Qq}{mL^3} x \;.
        \end{align*}
        %
        よって，
        %
        \begin{align*}
            \uwave{T = 2\pi \sqrt{\frac{mL^3}{2k_0 Qq}}} \;.
        \end{align*}

\end{enumerate}

\end{document}
